% MAIN SNAPSHOT — LSPV-5::1801::2026-10-02T18:01:15-07:00
\documentclass[11pt]{article}
\usepackage[margin=.86in]{geometry}
\usepackage{amsmath,amssymb}
\begin{document}
\section*{LSPV-5 MAIN Snapshot 1801}
\textbf{STATUS: OPEN.} P versus NP remains open. No unrestricted SAT lower bound, arbitrary-algorithm acquisition theorem, source-matched no-sharing theorem, nonrelativizing hinge, or nonalgebrizing hinge is established.

\section*{CORPUS RECONCILIATION}
\textbf{Corpus surveyed.} The current repository tree was enumerated at 345 entries. Governance files AGENTS.md, CODEX.md, and the pnp-research skill were fetched and obeyed. Root main.tex, TASK-001--007, verified MAIN history through MAIN 1658, all current LSPV-1--4 directories and materially relevant historical reports, HEC/VSC-HEC/G-PPC-HEC/CRS-PPC/CCSPC/RCCN, product-span and random-challenge explorations, cited experiment/result files, current GROK through 17:19, and the relevant CLAUDE review were surveyed. LSPV-4 LEGALFLAT1 remains the newest verified predecessor report; no failed LSPV-1--4 scientific write requiring reconstruction was found.

\textbf{New verified results since MAIN 1658.} No new LSPV-1--4 theorem report appeared after MAIN 1658. FULLPAIR6 remains PROVED/EXACT only in its finite declared source-typed scope, and RAMDOM1 remains PROVED/EXACT only as reconstruction accounting under its hypotheses. The tested ordered-DAG counts \(720,2776,5342,7771\) remain COMPUTATIONAL upper bounds, not global minima.

\textbf{Conflict resolved.} Root independent-review text still called MAIN 1158-R1 the verified arbiter and treated MAIN 1658 as non-arbiter. MAIN 1658 itself has the required MAIN header and verification sentinel and was created by LSPV-5. This run therefore adopts MAIN 1658 as the immediate predecessor snapshot. The scientific status is unchanged because both snapshots keep all load-bearing arrows open.

\textbf{Stale claims removed or downgraded.} The earlier \(2107\)-node \(D(4,8)\) construction remains valid for its declared order but is not the best tested ordered upper bound. The \(720\)-node result is still not a minimum FBDD, TDD, structured d-DNNF, or circuit size. Semantic mask counts are not machine-state or runtime lower bounds.

\textbf{Falsifications retained.} Generic authenticated causal replay, STORE-or-REACQUIRE alone, polynomial-transcript audit as a superpolynomial route, ordinary evolving restrictions as a universal transfer, explicit prime-product expansion, free canonicalization, Gram-to-characteristic-two signed lift, optional failed-shortcut costs, and raw incidence/tape geometry remain rejected or retired absent a new model-matched theorem.

\textbf{Open arrows changed.} None closed. The structural branch still needs a source-legal reconstruction-compatible characteristic-two growth theorem. The causal branch still needs forced SAME_REALIZATION acquisition against every compact evaluator and then a source-matched no-sharing theorem for arbitrary static SAT algorithms.

\textbf{Adopted handoffs.} FULLPAIR6 finite incidence recovery; MATRIXCERT1/EXCESS1/TURANEX1/COVERBREAK1 controls; ENVELOPE1/RAMDOM1 conservative reconstruction accounting; LEGALFLAT1/STRATFLAT1 flattening controls; TDD/structured-d-DNNF and arithmetic straight-line representations as mandatory hostile alternatives.

\textbf{Rejected handoffs.} Finite DAG growth, clause saturation, partition rank, LCD/permanent diagnostics, QBF strategy size, cryptographic pointer chasing, decision-tree direct products, or bounded adaptive-query lower bounds are not transferred to unrestricted SAT without an exact compiler.

\section*{Fresh literature audit}
Current primary literature continues to sharpen model boundaries rather than close the project hinge. Agarwal--Ben-David--Blais (arXiv:2609.12164) prove an information/amortized-size equivalence and a perfect direct-product theorem for a success-conditioned randomized decision-tree measure. This is a query-model theorem, not an unrestricted SAT no-sharing theorem. Capelli et al. (arXiv:2604.05537, SAT 2026) place TDDs between OBDD-style and structured deterministic DNNF representations and obtain FPT-size TDDs for bounded-treewidth CNFs. Berkholz--Micun--Razgon (arXiv:2609.21842) give quasipolynomial TDD-to-OBDD simulation. Ren--Williams (ECCC TR26-118) obtain near-maximum circuit lower bounds in a class with promise-MA access and analyze bounded rounds of adaptive NP queries with bounded witness lengths; this does not force arbitrary SAT algorithms into that model. Pebbling/proof-complexity and cryptographic sequential-memory results remain model-specific and cannot supply unconditional SAT acquisition. Recent algebraic-branching-program work likewise emphasizes that even superpolynomial lower bounds for multilinear ABPs remain open, so partition-rank diagnostics cannot be silently promoted.

\section*{Shortest surviving proof chain}
\[
\text{exact source-typed SU/vMS family}
\Longrightarrow
\operatorname{rank}_{\mathbb F_2}\Phi=\omega(m)
\Longrightarrow
\text{forced SAME\_REALIZATION acquisition}
\]
\[
\Longrightarrow
\text{source-matched no-sharing}
\Longrightarrow
\text{fully charged reconstruction}
\Longrightarrow
\text{nonrelativizing/nonalgebrizing contradiction}.
\]
Every load-bearing implication displayed above is OPEN or CONDITIONAL.

\section*{Structural branch}
FULLPAIR6 gives, for an independent legal direction set \(T=\{t_1,\ldots,t_r\}\), exact recovery of the complete legal pair-co-occurrence graph from legal slice matrices and the finite certificate
\[
\rho(H)\ge 1+r+|E_{\rm full}|.
\]
It does not prove \(r+|E_{\rm full}|=\omega(m)\). The next structural experiment must use one exact reconstruction-compatible SU/vMS orbit, emit the legal incidence data, and run COVERBREAK1 while searching for an intrinsically characteristic-two certificate such as an odd minor or binary dual obstruction.

\section*{Causal branch}
HEC/CRS-PPC/CCSPC remain useful only as a specification of what acquisition would have to force. LEGALFLAT1 shows that transcript-dependent legality does not evade flattening when legality, verification, and a polynomial-size strategy evaluator are polynomial time. Therefore a surviving causal theorem must exclude compact Skolem circuits, algebraic evaluators, TDDs, structured d-DNNFs, source/history quotients, arithmetic straight-line programs, memoization, and recomputation on the same source-matched family, then separately compile every hypothetical static SAT decider into the charged resource.

\section*{Conservative reconstruction ledger}
Under the declared ENVELOPE1/RAMCOMP1 hypotheses for one fixed SAT-decider exponent \(K\),
\[
D(K)=\max\{a,q+bK\},\qquad R(K)=\ell+c_0K^2,
\]
and RAMDOM1 states that the nested route is strictly cheaper exactly when
\[
\ell+c_0K^2<\max\{a,q+bK\}.
\]
Alternative routes are minimized, not added. \(K\) cannot be inflated or varied to manufacture a contradiction.

\section*{Single highest-value team target}
On one exact reconstruction-compatible SU/vMS polynomial-bit orbit, construct the source-transcribed binary operator \(\Phi\) and legal slice matrices and prove or falsify
\[
\boxed{\operatorname{rank}_{\mathbb F_2}\Phi=\omega(m)}
\]
with an intrinsically characteristic-two certificate, while simultaneously searching the same family for polynomial-size TDD/structured-d-DNNF, Skolem, algebraic, source/history, arithmetic, memoized, or recomputed representations. A compact representation falsifies that provenance route for the family. It is not progress toward \(P\neq NP\).

\section*{Barrier ledger}
Relativization remains open at the load-bearing acquisition/no-sharing step. Algebrization remains open; characteristic-two source dependence alone does not establish nonalgebrization. Natural-Proofs issues remain relevant to any constructive circuit property. Restricted decision-tree, branching-program, proof-complexity, communication, pebbling, online/streaming, and cryptographic lower bounds remain model-specific until a charged SAME_REALIZATION compiler is proved.

\section*{Final status}
P versus NP remains OPEN.

% END_OF_VERIFIED_MAIN::1801
\end{document}
